% ============================================================
% PARTIE III - CHAPITRE I : L'évaluation du projet
% ============================================================
\chapter{L'évaluation du projet}
\pagestyle{styleGlobal}
\label{chap:evaluation}

% ============================================================
\section{Organisation du projet}
% ============================================================

« SALISA » a été réalisé dans le cadre de l'obtention du diplôme de Licence Professionnelle
en Génie Logiciel à l'École Supérieure de Gestion et d'Administration des Entreprises (ESGAE)
de Brazzaville. Il s'agit d'un projet de fin d'études mené en binôme, sous la supervision d'un
Directeur de Projet (DP).

La conduite du projet repose sur des séances de travail hebdomadaires avec le Directeur de
Projet, tenues chaque samedi. Ces séances permettent de faire le point sur l'avancement, de
valider les livrables et d'orienter les prochaines étapes. La méthode de développement adoptée
est le \textbf{2TUP}, combinée au langage de modélisation \textbf{UML}.

% ============================================================
\section{Intervenants}
% ============================================================

Plusieurs acteurs ont participé à la réalisation de « SALISA ». Le tableau~\ref{tab:intervenants}
les présente avec leur fonction et leur rôle dans le projet.

\begin{table}[H]
	\centering
	\begin{tabular}{|p{5.5cm}|p{4cm}|p{4cm}|}
		\hline
		\rowcolor{vertCongo}
		\textcolor{white}{\textbf{Intervenant}} &
		\textcolor{white}{\textbf{Fonction}} &
		\textcolor{white}{\textbf{Titre / Rôle}} \\
		\hline
		M. Christophe Darly MASSAMBA BOUESSO & Étudiant en LPGL & Concepteur et développeur \\
		\hline
		M. Grâce Deo BATA & Étudiant en LPGL & Concepteur et développeur \\
		\hline
		M. Ornel Djenifer KIGOMA & Ingénieur informaticien & Directeur de projet, enseignant à l'ESGAE \\
		\hline
	\end{tabular}
	\caption{Intervenants du projet « SALISA »}
	\label{tab:intervenants}
\end{table}

% ============================================================
\section{Planification des tâches}
% ============================================================

La réalisation de « SALISA » a été découpée en plusieurs phases, couvrant la période allant
du 5 janvier 2026, date d'affichage officiel des binômes et de leurs thèmes, jusqu'à la soutenance
prévue le samedi 20 juin 2026. Le tableau~\ref{tab:planification} présente les grandes étapes
de ce planning.

\begin{table}[H]
	\centering
	\begin{tabular}{|p{7cm}|p{2.8cm}|p{2.8cm}|}
		\hline
		\rowcolor{vertCongo}
		\textcolor{white}{\textbf{Tâche}} &
		\textcolor{white}{\textbf{Début}} &
		\textcolor{white}{\textbf{Fin}} \\
		\hline
		Cadrage du projet et du thème & 05/01/2026 & 14/02/2026 \\
		\hline
		Étude de l'existant et analyse des besoins & 14/02/2026 & 15/03/2026 \\
		\hline
		Maquette et prototype de l'interface & 20/02/2026 & 20/03/2026 \\
		\hline
		Modélisation UML (diagrammes) & 15/03/2026 & 30/04/2026 \\
		\hline
		Modélisation de la base de données & 01/04/2026 & 15/04/2026 \\
		\hline
		Rédaction du rapport de soutenance & 20/02/2026 & 05/06/2026 \\
		\hline
		Développement back-end & \multicolumn{2}{c|}{\textit{à renseigner}} \\
		\hline
		Développement front-end & \multicolumn{2}{c|}{\textit{à renseigner}} \\
		\hline
		Intégration et déploiement & \multicolumn{2}{c|}{\textit{à renseigner}} \\
		\hline
		Tests et corrections & \multicolumn{2}{c|}{\textit{à renseigner}} \\
		\hline
		Préparation de la présentation & 05/06/2026 & 15/06/2026 \\
		\hline
		Révision finale et dépôt du dossier & \multicolumn{2}{c|}{avant le 02/06/2026} \\
		\hline
		Préparation et répétitions de soutenance & 10/06/2026 & 19/06/2026 \\
		\hline
		\rowcolor{vertClair}
		\textbf{Soutenance} & \textbf{20/06/2026} & \textbf{20/06/2026} \\
		\hline
	\end{tabular}
	\caption{Planification des tâches du projet « SALISA »}
	\label{tab:planification}
\end{table}

% ============================================================
\section{Diagramme de Gantt}
% ============================================================

\textit{[À insérer : diagramme de Gantt représentant le planning détaillé du projet, réalisé avec MS Project]}

% ============================================================
\section{Estimation des charges}
% ============================================================

L'estimation des charges regroupe l'ensemble des ressources nécessaires à la réalisation de
« SALISA ». Elle prend en compte les ressources humaines, les services en ligne, les logiciels
et les équipements matériels. Les montants sont exprimés en francs CFA.

\begin{table}[H]
	\centering
	\begin{tabular}{|p{4.5cm}|c|p{4.5cm}|p{3.2cm}|}
		\hline
		\rowcolor{vertCongo}
		\textcolor{white}{\textbf{Ressources}} &
		\textcolor{white}{\textbf{Quantité / Mois}} &
		\textcolor{white}{\textbf{Prix unitaire (en FCFA)}} &
		\textcolor{white}{\textbf{Prix total (en FCFA)}} \\
		\hline
		\rowcolor{grisclair}
		\multicolumn{4}{|l|}{\textbf{Ressources humaines}} \\
		\hline
		2 Développeurs & 5 & \raggedleft 300 000 & \raggedleft 3 000 000 \tabularnewline
		\hline
		\rowcolor{grisclair}
		\multicolumn{4}{|l|}{\textbf{Services en ligne}} \\
		\hline
		Connexion internet & 12 & \raggedleft 25 000 & \raggedleft 300 000 \tabularnewline
		\hline
		Serveur VPS (hébergement) & 12 & \raggedleft 15 000 & \raggedleft 180 000 \tabularnewline
		\hline
		\rowcolor{grisclair}
		\multicolumn{4}{|l|}{\textbf{Ressources logicielles}} \\
		\hline
		LaTeX / TeXstudio & 1 & \raggedleft Gratuit & \raggedleft Gratuit \tabularnewline
		\hline
		VS Code & 1 & \raggedleft Gratuit & \raggedleft Gratuit \tabularnewline
		\hline
		PostgreSQL & 1 & \raggedleft Gratuit & \raggedleft Gratuit \tabularnewline
		\hline
		Node.js / Redis & 1 & \raggedleft Gratuit & \raggedleft Gratuit \tabularnewline
		\hline
		Git et GitHub & 1 & \raggedleft Gratuit & \raggedleft Gratuit \tabularnewline
		\hline
		Docker & 1 & \raggedleft Gratuit & \raggedleft Gratuit \tabularnewline
		\hline
		Draw.io Desktop & 1 & \raggedleft Gratuit & \raggedleft Gratuit \tabularnewline
		\hline
		Postman (tests API) & 1 & \raggedleft Gratuit & \raggedleft Gratuit \tabularnewline
		\hline
		PlantUML & 1 & \raggedleft Gratuit & \raggedleft Gratuit \tabularnewline
		\hline
		MS Project (abonnement annuel Plan 1) & 12 & \raggedleft 6 250 & \raggedleft 75 000 \tabularnewline
		\hline
		\rowcolor{grisclair}
		\multicolumn{4}{|l|}{\textbf{Ressources matérielles}} \\
		\hline
		Ordinateur portable & 2 & \raggedleft 650 000 & \raggedleft 1 300 000 \tabularnewline
		\hline
		\rowcolor{vertCongo}
		\multicolumn{3}{|r|}{\textcolor{white}{\textbf{Montant total}}} &
		\multicolumn{1}{r|}{\textcolor{white}{\textbf{4 855 000}}} \\
		\hline
	\end{tabular}
	\caption{Estimation des charges du projet « SALISA »}
	\label{tab:estimation_charges}
\end{table}